\makeabstract
    {
    Synthesizing Naphthalenediimide-Based Self-Assembling Nanotubes
    }
    {
    Miraldy Luna\textsuperscript{1}, Antwone Pollard\textsuperscript{1}, Suntali Donahue\textsuperscript{1}, Professor David E. Hansen\textsuperscript{2}
    }
    {
    \textsuperscript{1} Amherst College, Amherst, MA\\
\textsuperscript{2} Department of Biochemistry, Amherst College, Amherst, MA
    }
    {
    With an emphasis on supramolecular chemistry, this summer we further investigated the serendipitous discovery in Sanders{\textquoteright} lab: amino acid NDI (Napthelenediimide) derivatives that form chiral helical nanotubes through reversible noncovalent interactions. Post DCM (dichloromethane) ban, our main focus has been finding an alternative solvent system to continue experimenting with different linkers for dimer and trimer formation. These iterations of NDI linkages have been shown to increase nanotube stability in previous Hansen lab research.
    }
    {
    Organic chemical synthesis typically involves chemical reactions, purification, and analysis. We utilized microwave chemistry to selectively add L-phenylalanine-t-butyl ester and L-serine-t-butyl ester to the NDI core following the Sander{\textquoteright}s lab technique. We also used flash chromatography to isolate the target compounds and NMR spectroscopy to confirm the synthesis of the target product.
    }
    {
    After testing various solvents, including 1,4-dioxane and acetonitrile, the latter appeared to be the most effective. We continued incorporating different solvophobic linkers to the proposed dimer. Providentially, we also discovered that using a 10-fold excess of diacyl chloride linker led to a more efficient monocoupling of the monomer control molecule, bypassing the previous need for complex organometallic chemistry in this synthesis.
    }
    {
    Future thesis work will focus on the trimer synthesis, utilizing similar methods and analysis. Alternative solvents would also need to be considered for the final steps of the synthesis, as part of a potential investigation into linker stability for nanotube formation. Given the ease of synthesizing and customizing NDI derivatives via microwave chemistry, this technique is very promising for continued research on synthetic self-assembling nanotubes for potential applications in drug delivery systems and catalytic reactions.
    }

    \newpage
    